% ═══════════════════════════════════════════════════════════════════════
% STREAM — PEARC 2026 Poster Extended Abstract (4 pages max)
% ═══════════════════════════════════════════════════════════════════════
% Deadline: March 23, 2026
% Submit via EasyChair: https://easychair.org/conferences/?conf=pearc26
% Same format as the short paper: \documentclass[manuscript]{acmart}
% ═══════════════════════════════════════════════════════════════════════

% PEARC26 submission format: single-column "manuscript" mode
% "review" adds line numbers for reviewer convenience
\documentclass[manuscript, review]{acmart}

\usepackage{booktabs}
\usepackage{graphicx}

\begin{document}

% ═══════════════════════════════════════════════════════════════════════
% TITLE & AUTHORS
% ═══════════════════════════════════════════════════════════════════════

\title{STREAM: Multi-Tier LLM Inference Middleware with Dual-Channel HPC Token Streaming}

\author{Anas Nassar}
\authornote{Corresponding author.}
\affiliation{
  \department{Advanced Cyberinfrastructure for Education and Research (ACER)}
  \institution{University of Illinois Chicago}
  \city{Chicago}
  \state{Illinois}
  \country{USA}
}
\email{nassar@uic.edu}

\author{Steve Mohr}
\affiliation{
  \department{Advanced Cyberinfrastructure for Education and Research (ACER)}
  \institution{University of Illinois Chicago}
  \city{Chicago}
  \state{Illinois}
  \country{USA}
}
\email{smohr@uic.edu}

\author{Himanshu Sharma}
\affiliation{
  \department{Advanced Cyberinfrastructure for Education and Research (ACER)}
  \institution{University of Illinois Chicago}
  \city{Chicago}
  \state{Illinois}
  \country{USA}
}
\email{himanshu@uic.edu}

% ═══════════════════════════════════════════════════════════════════════
% ABSTRACT
% ═══════════════════════════════════════════════════════════════════════

\begin{abstract}
We present STREAM, a middleware that routes user queries across three LLM inference tiers---local (Ollama), campus HPC (vLLM via Globus Compute), and cloud (commercial APIs)---based on query complexity. STREAM introduces a novel dual-channel architecture that separates Globus Compute's batch-oriented control plane from a WebSocket-based data plane, enabling real-time token streaming from HPC resources. This poster demonstrates the system architecture, the streaming relay design, and preliminary latency measurements comparing relay-based streaming against Globus Compute's default batch mode.
\end{abstract}

\keywords{LLM inference, HPC, Globus Compute, token streaming, WebSocket, tiered routing}

\maketitle

% ═══════════════════════════════════════════════════════════════════════
% 1. INTRODUCTION & MOTIVATION
% ═══════════════════════════════════════════════════════════════════════

\section{Introduction}

Campus HPC clusters increasingly host powerful GPUs capable of running large open-source LLMs, but their batch-oriented job schedulers create a poor user experience for interactive chat. Globus Compute~\cite{globuscompute2024} provides federated authentication and remote function execution on HPC resources, but its execution model requires waiting for complete results---there is no mechanism for streaming partial outputs.

FIRST~\cite{first2025}, the closest prior work, uses Globus Compute for HPC-based LLM inference at Argonne National Laboratory but reports a median response latency of 16.3 seconds, with users waiting for the entire response before seeing any output. Campus deployments at Dartmouth~\cite{dartmouth2025} and Purdue~\cite{purdue2025} avoid this problem by running inference servers directly accessible over HTTP, but sacrifice federated authentication and multi-institutional access.

STREAM (Smart Tiered Routing Engine for AI Models) solves the HPC streaming problem while preserving Globus Compute's authentication and access control benefits. It routes queries across three tiers---local device, campus HPC, and cloud---with automatic complexity-based routing, and introduces a dual-channel architecture that enables real-time token streaming from HPC resources.

% ═══════════════════════════════════════════════════════════════════════
% 2. SYSTEM ARCHITECTURE
% ═══════════════════════════════════════════════════════════════════════

\section{System Architecture}

STREAM classifies each query as LOW, MEDIUM, or HIGH complexity using a local LLM-as-judge with keyword-based fallback. LOW queries route to local Ollama (free, private); MEDIUM to the campus HPC cluster via Globus Compute; HIGH to cloud APIs. If a tier is unavailable, the router falls back to the next tier. Table~\ref{tab:tiers} summarizes the tiers.

\begin{table}[h]
\centering
\caption{STREAM inference tiers.}
\label{tab:tiers}
\small
\begin{tabular}{@{}lllr@{}}
\toprule
\textbf{Tier} & \textbf{Model} & \textbf{Hardware} & \textbf{Cost} \\
\midrule
Local     & Llama 3.2 3B / Gemma 3 4B & User's CPU    & \$0 \\
Lakeshore & Qwen2.5-VL 72B            & H100 96GB GPU & \$0 \\
Cloud     & Claude Sonnet / GPT-4      & Commercial API & \$\$\$ \\
\bottomrule
\end{tabular}
\end{table}

STREAM supports two deployment modes sharing 95\% of the codebase: a \textbf{desktop mode} (standalone application via PyWebView where the local tier runs on the user's own machine) and a \textbf{server mode} (Docker Compose serving multiple users). Both modes use the same routing, streaming, and dual-channel architecture.

% ═══════════════════════════════════════════════════════════════════════
% 3. DUAL-CHANNEL HPC STREAMING
% ═══════════════════════════════════════════════════════════════════════

\section{Dual-Channel HPC Streaming}

The key contribution of this work is a dual-channel architecture that enables real-time token streaming through Globus Compute without modifying the Globus platform itself.

\textbf{Control plane} (Globus Compute): Handles OAuth2 authentication, job submission via AMQP message queues, and status reporting. The middleware submits the inference function to Globus Compute, which dispatches it to the HPC endpoint.

\textbf{Data plane} (WebSocket relay): A lightweight relay server that forwards tokens in real-time. The remote function on the HPC node connects to the relay as a \textit{producer}, streaming tokens from vLLM as they are generated. The middleware connects as a \textit{consumer}, receiving tokens and forwarding them to the user's browser as Server-Sent Events (SSE).

% TODO: Replace with Excalidraw diagram for the poster
% \begin{figure}[h]
%     \centering
%     \includegraphics[width=\textwidth]{dual-channel.png}
%     \caption{Dual-channel architecture: control plane (Globus Compute) for auth and dispatch, data plane (WebSocket relay) for real-time token streaming.}
%     \label{fig:dualchannel}
% \end{figure}

Both the producer (on the HPC node behind a firewall) and the consumer connect \textit{outbound} to the relay, eliminating the need for inbound firewall rules on HPC infrastructure. Each request uses a unique channel ID for isolation. If the relay is unavailable, STREAM falls back to batch mode automatically.

% ═══════════════════════════════════════════════════════════════════════
% 4. PRELIMINARY RESULTS & POSTER CONTENT
% ═══════════════════════════════════════════════════════════════════════

\section{Preliminary Results}

% TODO: Fill in with actual measurements
Table~\ref{tab:latency} shows preliminary latency measurements comparing relay-based streaming against Globus Compute's default batch mode for the Lakeshore HPC tier.

\begin{table}[h]
\centering
\caption{Lakeshore HPC tier latency: relay streaming vs.\ batch mode.}
\label{tab:latency}
\small
\begin{tabular}{@{}lrr@{}}
\toprule
\textbf{Mode} & \textbf{TTFT (s)} & \textbf{Total (s)} \\
\midrule
WebSocket relay (streaming) & TODO & TODO \\
Batch mode (no streaming)   & TODO & TODO \\
\bottomrule
\end{tabular}
\end{table}

The poster will present: (1)~the full system architecture diagram showing all three tiers and the dual-channel design, (2)~a detailed sequence diagram of the relay handshake and token flow, (3)~latency comparisons between relay streaming and batch mode, and (4)~a video demonstration accessible via the QR code in Figure~\ref{fig:qr}.

% TODO: Generate a QR code image from your YouTube demo link.
%   1. Go to https://www.qrcode-monkey.com/ (or any QR generator)
%   2. Paste your YouTube URL
%   3. Download as PNG
%   4. Upload to Overleaf
%   5. Uncomment the includegraphics line below
\begin{figure}[h]
    \centering
    % \includegraphics[width=0.25\textwidth]{qr-demo.png}
    \fbox{\parbox{0.2\textwidth}{\centering\vspace{1.5cm}
        \textit{QR code}\\
        \textit{(demo video)}
    \vspace{1.5cm}}}
    \caption{Scan to watch STREAM demo: real-time token streaming from UIC's Lakeshore HPC cluster via the dual-channel architecture.}
    \label{fig:qr}
\end{figure}

\begin{acks}
The authors thank Leonard Apanasevich, Assistant Director of Research Facilitation at ACER, for his support. This research used resources of the Lakeshore HPC cluster at the University of Illinois Chicago.
\end{acks}

\subsection*{AI Tool Disclosure}
Claude Code (Anthropic's CLI-based AI coding assistant) was used iteratively during STREAM's development for implementation assistance and code review. All system design decisions, architectural innovations, and evaluation are the original work of the authors.

% ═══════════════════════════════════════════════════════════════════════
% REFERENCES
% ═══════════════════════════════════════════════════════════════════════

\begin{thebibliography}{9}

\bibitem{first2025}
K.~Hippe et al.,
``FIRST: Federated Inference Resource Scheduling Toolkit for Scientific AI Model Access,''
in \textit{Proc. SC'25 Workshops}, ACM, 2025.

\bibitem{dartmouth2025}
S.~Bratman et al.,
``Dartmouth Chat---Deploying an Open-Source LLM Stack at Scale,''
in \textit{Proc. PEARC'25}, ACM, 2025.

\bibitem{purdue2025}
X.~Wang et al.,
``Providing On-Prem GenAI Inference Services to a Campus Community,''
in \textit{Proc. PEARC'25}, ACM, 2025.

\bibitem{globuscompute2024}
R.~Chard et al.,
``Globus Compute: A Federated Function-as-a-Service Platform for Computational Science,''
\textit{Future Generation Computer Systems}, vol.~154, pp.~558--572, 2024.

\bibitem{routellm2024}
I.~Ong et al.,
``RouteLLM: Learning to Route LLMs with Preference Data,''
in \textit{Proc. ICLR}, 2025.

\end{thebibliography}

\end{document}
